% ch02_u01: 原书 p.15--p.23
\SourcePage{15}

\bdsec{2.1. $\operatorname{Bun}_G$ 的几何}

我们沿用 0.1 的记号；特别地，$G$ 是半单群，$X$ 是亏格 $g>1$ 的光滑射影曲线。

\bdp{2.1.1} 众所周知，$\operatorname{Bun}_G$ 是纯维数 $(g-1)\dim G$ 的光滑代数栈。$\operatorname{Bun}_G$ 的连通分支集合可（经「第一陈类」映射）典范地等同于 $H^{2}(X,\pi_{1}^{et}(G))=\pi_{1}(G)$。这里 $\pi_{1}^{et}(G)$ 是 Grothendieck 意义下的基本群；$\pi_{1}(G)$ 是 $G$ 的余权群模去余根子群所得的商。二者相差一个 Tate 扭：$\pi_{1}^{et}(G)=\pi_{1}(G)(1)$。

对 $\mathcal{F}\in\operatorname{Bun}_G$，切层 $\Theta=\Theta_{\operatorname{Bun}_G}$ 在 $\mathcal{F}$ 处的纤维是 $H^{1}(X,\mathfrak{g}_{\mathcal{F}})$。说明一下：对 $G$-模 $W$，记 $W_{\mathcal{F}}$ 为 $W$ 的 $\mathcal{F}$-扭，它是 $X$ 上的向量丛。这里经伴随作用把 $\mathfrak{g}$ 看作 $G$-模。

按定义，典范线丛 $\omega=\omega_{\operatorname{Bun}_G}$ 是 $\operatorname{Bun}_G$ 的余切复形的行列式（见 [LMB93]）。该复形在 $\mathcal{F}\in\operatorname{Bun}_G$ 上的纤维与 $R\Gamma(X,\mathfrak{g}_{\mathcal{F}})[1]$ 对偶（见 [LMB93]），故 $\omega$ 在 $\mathcal{F}$ 上的纤维是 $\det R\Gamma(X,\mathfrak{g}_{\mathcal{F}})$。\footnote{作者们不该忘记核验 [LMB93] 中确实包含此处所断言的内容！！}

\bdp{2.1.2} \emph{命题.} $\operatorname{Bun}_G$ 在 1.1.1 的意义下是极良的。

证明将在 2.10.5 给出。事实上，我们将用到 $\operatorname{Bun}_G$ 是良的这一事实。由 1.1，有泊松代数层 $P=P_{\operatorname{Bun}_G}$ 与扭微分算子层 $D^{\lambda}=D_{\operatorname{Bun}_G,\omega^{\lambda}}$。众所周知，当 $\lambda\neq 1/2$ 时，$D^{\lambda}$ 的整体截面只有局部常值函数。在第 2 节与第 3 节中我们处理 $D':=D^{1/2}$；其截面就简称为扭微分算子。

\bdsec{2.2. Hitchin 构造 I}

\SourcePage{16}

\bdp{2.2.1} 置 $\mathcal{C}=\mathcal{C}_{\mathfrak{g}}:=\operatorname{Spec}(\operatorname{Sym}\mathfrak{g})^{G}$；这是 $\mathfrak{g}^{*}$ 关于余伴随作用的仿射概形商。$\mathcal{C}$ 带有乘法群 $\mathbb{G}_m$ 的典范作用，它来自 $\mathfrak{g}^{*}$ 上的位似变换。非典范地取一组次数为 $d_i$、$i\in I$ 的齐次生成元 $p_i\in(\operatorname{Sym}\mathfrak{g})^{G}$，就把 $\mathcal{C}$ 与坐标空间 $\mathbb{C}^{I}$ 等同；元素 $\lambda\in\mathbb{G}_m$ 通过对角矩阵 $(\lambda^{d_i})$ 作用。

\bdp{2.2.2} 记 $\mathcal{C}_{\omega_X}$ 为 $\mathcal{C}$ 关于上述 $\mathbb{G}_m$-作用的 $\omega_X$-扭（这里把典范丛 $\omega_X$ 看作 $X$ 上的 $\mathbb{G}_m$-torsor，即主齐性空间）。它是 $X$ 上的丛；上述 $p_i$ 把 $\mathcal{C}_{\omega_X}$ 与 $\prod_I\omega_X^{\otimes d_i}$ 等同。置
\[
\operatorname{Hitch}(X)=\operatorname{Hitch}_{\mathfrak{g}}(X):=\Gamma(X,\mathcal{C}_{\omega_X}).
\]
换言之，$\operatorname{Hitch}(X)=\operatorname{Mor}((\operatorname{Sym}^{\cdot}\mathfrak{g})^{G},\Gamma(X,\omega_X^{\otimes\cdot}))$（分次代数的态射）。把 $\operatorname{Hitch}(X)$ 看作带 $\mathbb{G}_m$-作用的代数簇；它非典范地同构于向量空间 $\prod_I\Gamma(X,\omega_X^{\otimes d_i})$。存在唯一的点 $0\in\operatorname{Hitch}(X)$，它在 $\mathbb{G}_m$ 的作用下不动。记 $\mathfrak{z}^{cl}(X)=\mathfrak{z}_{\mathfrak{g}}^{cl}(X)$ 为 $\operatorname{Hitch}(X)$ 上的函数环；这是分次交换代数。更确切地说，$\mathfrak{z}^{cl}(X)$ 的分次对应于 $\mathfrak{z}^{cl}(X)$ 上一个 $\mathbb{G}_m$-作用，它与 $\mathcal{C}$ 上 $\mathbb{G}_m$-作用诱导的作用相反；因此 $\mathfrak{z}^{cl}(X)$ 的分次为正。

\bdp{2.2.3} 由 Serre 对偶与 2.1.1，$\mathcal{F}\in\operatorname{Bun}_G$ 处 $\operatorname{Bun}_G$ 的余切空间 $T_{\mathcal{F}}^{*}\operatorname{Bun}_G$ 等于 $\Gamma(X,\mathfrak{g}_{\mathcal{F}}^{*}\otimes\omega_X)$。$G$-不变投影 $\mathfrak{g}^{*}\to\mathcal{C}$ 给出态射 $\mathfrak{g}_{\mathcal{F}}^{*}\otimes\omega_X\to\mathcal{C}_{\omega_X}$，以及映射 $p_{\mathcal{F}}:T_{\mathcal{F}}^{*}\operatorname{Bun}_G\to\operatorname{Hitch}(X)$。当 $\mathcal{F}$ 变动时得到态射
\[
p:T^{*}\operatorname{Bun}_G\to\operatorname{Hitch}(X)
\]
或等价地，分次交换代数的态射
\[
h_X^{cl}:\mathfrak{z}^{cl}(X)\to\Gamma(T^{*}\operatorname{Bun}_G,\mathcal{O})=\Gamma(\operatorname{Bun}_G,P).
\]
$p$ 称为 Hitchin 纤维化。

\SourcePage{17}

记 $\operatorname{Bun}_G^{\gamma}$ 为 $\operatorname{Bun}_G$ 中对应于 $\gamma\in\pi_1(G)$ 的连通分支（见 2.1.1），记 $p^{\gamma}$ 为 $p$ 到 $T^{*}\operatorname{Bun}_G^{\gamma}$ 的限制。

\bdp{2.2.4} \emph{定理.}（[Hit87]，[Fal93]，[Gi97]）
\begin{enumerate}[label=(\roman*),leftmargin=2.6em]
\item $h_X^{cl}$ 的像由泊松交换的函数组成。
\item $\dim\operatorname{Hitch}(X)=\dim\operatorname{Bun}_G=(g-1)\cdot\dim\mathfrak{g}$。
\item $p$ 是平坦的，其纤维的维数纯为 $\dim\operatorname{Bun}_G$。对每个 $\gamma\in\pi_1(X)$，$p^{\gamma}$ 是满射。
% [Tutor Note: 原书 p.17 定理 2.2.4 的 (iii) 与 (v) 均印作 γ ∈ π₁(X)，按上下文应为 π₁(G)；底本如此，保留原样。]
\item 存在非空开集 $U\subset\operatorname{Hitch}(X)$，使得对任意 $\gamma\in\pi_1(G)$，态射 $(p^{\gamma})^{-1}(U)\to U$ 固有且光滑，并且其纤维连通。事实上，$p^{\gamma}$ 在 $u\in U$ 上的纤维同构于某个阿贝尔簇 $A_u$ 与中心 $Z\subset G$ 的分类栈之积。
\item 对每个 $\gamma\in\pi_1(X)$，态射 $\mathfrak{z}^{cl}(X)\to\Gamma(\operatorname{Bun}_G^{\gamma},P)$ 是同构。\hfill$\Box$
\end{enumerate}

\emph{注.}
\begin{enumerate}[label=(\roman*),leftmargin=2.6em]
\item 毋庸赘言，定理 2.2.4 的主要贡献属于 Hitchin [Hit87]。
\item 定理 2.2.4 蕴含 $p$ 是 Lagrange 纤维化；或者，若你愿意这么说，来自 $h_X^{cl}(\mathfrak{z}^{cl}(X))$ 的哈密顿量在 $T^{*}\operatorname{Bun}_G$ 上定义完全可积系统。在栈的语境中我们不怕用这些词，因为 Lagrange 纤维化的概念是双有理的。又因 $\operatorname{Bun}_G$ 在 1.1.1 的意义下是极良的，$T^{*}\operatorname{Bun}_G$ 有开稠密的 Deligne--Mumford 子栈 $T^{*}\operatorname{Bun}_G^{0}$，它在显然的意义下是辛的（这里 $\operatorname{Bun}_G^{0}$ 是自同构群有限的 $G$-丛构成的栈）。
\item Hitchin 在 [Hit87] 中给出陈述 (i) 的复分析证明。我们将在 2.4.3 给出 (i) 的代数证明。
\SourcePage{18}
\item Hitchin 对 (ii) 的证明很简单：由 2.2.2，$\dim\operatorname{Hitch}(X)=\sum_i\dim\Gamma(X,\omega_X^{\otimes d_i})$；因 $g>1$，$\dim\Gamma(X,\omega_X^{\otimes d_i})=(g-1)(2d_i-1)$；最后 $(g-1)\sum_i(2d_i-1)=(g-1)\dim\mathfrak{g}=\dim\operatorname{Bun}_G$。
\item 对经典群 $G$，陈述 (iv) 已由 Hitchin [Hit87] 证明。一般情形由 Faltings 证明（[Fal93] 中的定理 III.2）。
\item 陈述 (v) 由 (iii) 与 (iv) 推出。
\item (iii) 的证明的一些评述将在 2.10 给出。
\end{enumerate}

\bdp{2.2.5} 我们的目标是解决如下量子化问题：构造滤交换代数 $\mathfrak{z}(X)$，并赋予它同构 $\sigma_{\mathfrak{z}(X)}:\operatorname{gr}\mathfrak{z}(X)\xrightarrow{\sim}\mathfrak{z}^{cl}(X)$，以及滤代数的态射 $h_X:\mathfrak{z}(X)\to\Gamma(\operatorname{Bun}_G,D')$；后者与象征映射相容，即 $\sigma_{\operatorname{Bun}_G}\circ\operatorname{gr}h_X=h_X^{cl}\circ\sigma_{\mathfrak{z}(X)}$（$\sigma_{\operatorname{Bun}_G}$ 的定义见 1.1.4 与 1.1.6）。

注意，2.2.4(v) 于是蕴含：对任意 $\gamma\in\pi_1(X)$，映射 $h_X^{\gamma}:\mathfrak{z}(X)\to\Gamma(\operatorname{Bun}_G^{\gamma},D')$ 是同构。
% [Tutor Note: 原书 p.18 此处亦印作 γ ∈ π₁(X)，按上下文应为 π₁(G)；底本如此，保留原样。]
因此若 $G$ 单连通，则这样的构造唯一；它归结为两条断言：$\Gamma(\operatorname{Bun}_G,D')$ 是交换代数；$T^{*}\operatorname{Bun}_G$ 上的任一整体函数都是某个整体扭微分算子的象征。

我们不知道如何用整体考虑直接解决此问题。我们将从 Hitchin 图景的局部版本出发，遵循 1.2 的量子化方案。上述量子化问题的同一个解将给出两个构造。第一个（见 2.5.5）较易表述；第二个（见 2.7.4）的优点是完全典范。为证明第一个构造确实给出解，我们使用第二个。正是第二个构造将给出 $\operatorname{Spec}\mathfrak{z}(X)$ 与 $X$ 上 $(^{L}G)_{ad}$-局部系统栈的某个子空间的等同（见 3.3.2）。

\bdsec{2.3. $\operatorname{Bun}_G$ 的几何 II}

让我们回忆 $\operatorname{Bun}_G$ 如何纳入 1.2.6 的框架。

\SourcePage{19}

\bdp{2.3.1} 固定一点 $x\in X$。记 $\mathcal{O}_x$ 为 $x$ 的完备局部环，$K_x$ 为其分式域。设 $\mathfrak{m}_x\subset\mathcal{O}_X$ 为极大理想。置 $\mathcal{O}_x^{(n)}:=\mathcal{O}_X/\mathfrak{m}_x^{n}$（于是 $\mathcal{O}_x=\varprojlim\mathcal{O}_x^{(n)}$）。群 $G(\mathcal{O}_x^{(n)})$ 是某个仿射代数群的 $\mathbb{C}$-点群；我们仍把它记作 $G(\mathcal{O}_x^{(n)})$（滥用记号）。$G(\mathcal{O}_x^{(n)})$ 是 $G(\mathcal{O}_x^{(n+1)})$ 的商。于是 $G(\mathcal{O}_x)=\varprojlim G(\mathcal{O}_x^{(n)})$ 是仿射群概形。

记 $\operatorname{Bun}_{G,nx}$ 为 $X$ 上在 $\operatorname{Spec}\mathcal{O}_x^{(n)}$ 上平凡化的 $G$-丛构成的栈（注意除子 $nx$ 与子概形 $\operatorname{Spec}\mathcal{O}_x^{(n)}\subset X$ 相同）。这是 $\operatorname{Bun}_G$ 上的 $G(\mathcal{O}_x^{(n)})$-torsor。把 $\operatorname{Bun}_{G,nx}$ 的点记作 $(\mathcal{F},\alpha^{(n)})$。有显然的仿射投影 $\operatorname{Bun}_{G,(n+1)x}\to\operatorname{Bun}_{G,nx}$。置 $\operatorname{Bun}_{G,\underline{x}}:=\varprojlim\operatorname{Bun}_{G,nx}$；这是 $\operatorname{Bun}_G$ 上的 $G(\mathcal{O}_x)$-torsor。

\bdp{2.3.2} \emph{命题.} $\operatorname{Bun}_{G,\underline{x}}$ 是概形。$\operatorname{Bun}_{G,\underline{x}}$ 上的 $G(\mathcal{O}_x)$-作用满足 1.2.6 的条件 (16)。\hfill$\Box$

\bdp{2.3.3} 众所周知，$\operatorname{Bun}_{G,\underline{x}}$ 上的 $G(\mathcal{O}_x)$-作用可典范地延拓为群归纳概形 $G(K_x)$ 的作用（归纳概形的定义见 7.11.1，归纳概形 $G(K_x)$ 的定义见 7.11.2 (iv)）。因 $\operatorname{Lie}G(K_x)=\mathfrak{g}\otimes K_x$，特别地有 Harish-Chandra 对 $(\mathfrak{g}\otimes K_x,G(\mathcal{O}_x))$ 在 $\operatorname{Bun}_{G,\underline{x}}$ 上的作用。

让我们回忆 $G(K_x)$-作用的定义。由 7.11.2 (iv)，须对任意 $\mathbb{C}$-代数 $R$ 定义 $G(R\widehat{\otimes}K_x)$ 在 $\operatorname{Bun}_{G,\underline{x}}(R)$ 上的作用。为此使用下述定理，它本质上属于 A. Beauville 与 Y. Laszlo。置 $X':=X\setminus\{x\}$。

\bdp{2.3.4} \emph{定理.} $X\otimes R$ 上的 $G$-丛 $\mathcal{F}$ 等同于三元组 $(\mathcal{F}_1,\mathcal{F}_2,\varphi)$：其中 $\mathcal{F}_1$ 是 $X'\otimes R$ 上的 $G$-丛，$\mathcal{F}_2$ 是 $\operatorname{Spec}(R\widehat{\otimes}\mathcal{O}_x)$ 上的 $G$-丛，$\varphi$ 是 $\mathcal{F}_1$ 与 $\mathcal{F}_2$ 到 $\operatorname{Spec}(R\widehat{\otimes}K_x)$ 的拉回之间的同构。

更确切地说，从 $X\otimes R$ 上 $G$-丛 $\mathcal{F}$ 的范畴（=群胚）到上述三元组 $(\mathcal{F}_1,\mathcal{F}_2,\varphi)$ 的范畴的函子，由 $\mathcal{F}_1:=\mathcal{F}|_{X'\otimes R}$、$\mathcal{F}_2:=\mathcal{F}$ 到 $\operatorname{Spec}(R\widehat{\otimes}\mathcal{O}_x)$ 的拉回、$\varphi:=\mathrm{id}$ 定义；它是范畴等价。

\SourcePage{20}

由该定理，$\operatorname{Bun}_{G,\underline{x}}$ 的 $R$-点等同于 $X'\otimes R$ 上的一个 $G$-丛，连同一个其到 $\operatorname{Spec}(R\widehat{\otimes}K_x)$ 的拉回的平凡化。于是 $G(R\widehat{\otimes}K_x)$ 通过改变平凡化作用在 $\operatorname{Bun}_{G,\underline{x}}(R)$ 上。这样就得到 $G(K_x)$ 在 $\operatorname{Bun}_{G,\underline{x}}$ 上的作用。

定理 2.3.4 的证明基于下述定理，它是 [BLa95] 主要结果的一个特例。

\bdp{2.3.5} \emph{定理.}（Beauville--Laszlo）平坦拟凝聚 $\mathcal{O}_{X\otimes R}$-模 $\mathcal{M}$ 的范畴等价于三元组 $(\mathcal{M}_1,\mathcal{M}_2,\varphi)$ 的范畴。其中 $\mathcal{M}_1$ 是 $X'\otimes R$ 上的平坦拟凝聚 $\mathcal{O}$-模，$\mathcal{M}_2$ 是 $\operatorname{Spec}(R\widehat{\otimes}\mathcal{O}_x)$ 上的平坦拟凝聚 $\mathcal{O}$-模，$\varphi$ 是 $\mathcal{M}_1$ 与 $\mathcal{M}_2$ 到 $\operatorname{Spec}(R\widehat{\otimes}K_x)$ 的拉回之间的同构（从第一个范畴到第二个范畴的函子按定理 2.3.4 定义）。$\mathcal{M}$ 局部自由且秩有限，当且仅当相应的 $\mathcal{M}_1$ 与 $\mathcal{M}_2$ 具有此性质。

\emph{注.} 若 $R$ 诺特且层凝聚，则有 M. Artin 的一个远为一般的「粘合定理」（[Ar] 中的定理 2.6）。但因 $G(K_x)$ 的子概形通常是无限型的，我们使用 Beauville--Laszlo 定理，它无需诺特假设。

为从 2.3.5 推出定理 2.3.4，只需把 $G$-丛解释为张量函子 $\{G\text{-模}\}\to\{\text{向量丛}\}$。或者，可把 $X\otimes R$ 上的 $G$-丛解释为主 $G$-丛，即平坦仿射态射 $\pi:\mathcal{F}\to X\otimes R$，连同 $G$ 在 $\mathcal{F}$ 上满足某些性质的作用；于是可用层 $\mathcal{M}:=\pi_{*}\mathcal{O}_{\mathcal{F}}$ 改写这些数据，并应用定理 2.3.5。
% [Tutor Note: 原书 p.20 此处印作 "principle G-bundle"，应为 "principal G-bundle"；底本如此，译文按正字处理。]

\bdp{2.3.6} \emph{注.} 这里给出 $\mathfrak{g}\otimes K_x$ 在 $\operatorname{Bun}_{G,\underline{x}}$ 上作用的直接描述，该作用由 $G(K_x)$ 的作用诱导（将来不会用到 ???）。取 $(\mathcal{F},\bar\alpha)\in\operatorname{Bun}_{G,\underline{x}}$，$\bar\alpha=\varprojlim\alpha^{(n)}$。$\operatorname{Bun}_{G,nx}$ 在 $(\mathcal{F},\alpha^{(n)})$ 处的切空间是 $H^{1}(X,\mathfrak{g}_{\mathcal{F}}(-nx))$，故 $\Theta_{\operatorname{Bun}_{G,\underline{x}}}$ 在 $(\mathcal{F},\bar\alpha)$ 处的纤维等于 $\varprojlim H^{1}(X,\mathfrak{g}_{\mathcal{F}}(-nx))=H^{1}_{c}(X\setminus\{x\},\mathfrak{g}_{\mathcal{F}})$。有通常的满射 $\mathfrak{g}_{\mathcal{F}}\otimes_{\mathcal{O}_X}K_x\to H^{1}_{c}(X\setminus\{x\},\mathfrak{g}_{\mathcal{F}})$。用 $\bar\alpha$ 把 $\mathfrak{g}_{\mathcal{F}}\otimes_{\mathcal{O}_X}K_x$ 与 $\mathfrak{g}\otimes K_x$ 等同。
% [Tutor Note: 原书 p.20 此处交叉引用印作「???」，是 LaTeX 未解析的引用；底本如此，保留原样。]

\SourcePage{21}

当 $(\mathcal{F},\bar\alpha)$ 变动时得到映射 $\mathfrak{g}\otimes K_x\to\Theta_{\operatorname{Bun}_{G,\underline{x}}}$。我们的 $\mathfrak{g}\otimes K_x$-作用是该映射的相反数（???）。
% [Tutor Note: 原书 p.21 此处交叉引用印作「???」；底本如此，保留原样。]

\bdp{2.3.7} \emph{注.} 设 $D\subset X\otimes R$ 是 $\operatorname{Spec}R$ 上的有限闭子概形，它可局部由一个方程定义（即 $D$ 是有效相对 Cartier 除子）。记 $\tilde{D}$ 为 $D$ 的形式邻域，$A$ 为 $\tilde{D}$ 的坐标环（于是 $\tilde{D}$ 是仿射形式概形，$\operatorname{Spec}A$ 是真概形）。则把 $X'\otimes R$ 换成 $(X\otimes R)\setminus D$、把 $R\widehat{\otimes}\mathcal{O}_x$ 换成 $A$、把 $\operatorname{Spec}(R\widehat{\otimes}K_x)$ 换成 $(\operatorname{Spec}A)\setminus D$ 后，定理 2.3.4 与 2.3.5 仍然成立。

若 $D$ 的正规丛平凡，这由 [BLa95] 的主要定理推出：事实上，此时可构造仿射邻域 $U\supset D$，使得在 $U$ 内子概形 $D$ 由整体方程 $f=0$ 定义，$f\in H^{0}(U,\mathcal{O}_U)$（这是 [BLa95] 中考虑的情形）。\footnote{为构造 $U$ 与 $f$，注意当 $n$ 充分大时存在 $\varphi_n\in H^{0}(X\otimes R,\mathcal{O}_{X\otimes R}(nD))$，使得 $\mathcal{O}_{X\otimes R}(nD)/\mathcal{O}_{X\otimes R}((n-1)D)$ 由 $\varphi_n$ 生成；然后置 $U:=(X\otimes R)\setminus\{\varphi_n\varphi_{n+1}\text{ 的零点集}\}$，$f:=\varphi_n/\varphi_{n+1}$（此构造在 $D\to\operatorname{Spec}R$ 为满射时有效，这是一个无害的假设）。} 对本工作的目的而言，$D$ 的正规丛平凡的情形已足够。为处理一般情形，需要 [BLa95] 主要定理的整体化版本（见 2.12）。此外还需把态射 $\tilde{D}\to X\otimes R$ 延拓为态射 $\operatorname{Spec}A\to X\otimes R$（显然这样的延拓唯一，但其存在性须证明）；见 2.12。

\bdsec{2.4. Hitchin 构造 II}

\bdp{2.4.1} 置 $\omega_{\mathcal{O}_x}:=\varprojlim\omega_{\mathcal{O}_x^{(n)}}$，其中 $\omega_{\mathcal{O}_x^{(n)}}$ 是 $\mathcal{O}_x^{(n)}=\mathcal{O}_x/\mathfrak{m}_x^{n}$ 的微分模。记 $\operatorname{Hitch}_x^{(n)}$ 为 $\mathcal{C}_{\omega_X}$ 在 $\operatorname{Spec}\mathcal{O}_x^{(n)}$ 上的截面概形。这是带 $\mathbb{G}_m$-作用的仿射概形，非典范地同构于向量空间 $M/\mathfrak{m}_x^{n}M$，其中 $M:=\prod\omega_{\mathcal{O}_x}^{\otimes d_i}$。置
\[
\operatorname{Hitch}_x=\operatorname{Hitch}_{\mathfrak{g}}(\mathcal{O}_x):=\varprojlim\operatorname{Hitch}_x^{(n)}.
\]
这是带 $\mathbb{G}_m$-作用的仿射概形，非典范地同构于 $M=\prod\omega_{\mathcal{O}_x}^{\otimes d_i}$。于是 $\operatorname{Hitch}_x$ 是 $\mathcal{C}_{\omega_X}$ 在 $\operatorname{Spec}\mathcal{O}_x$ 上的截面概形。

\SourcePage{22}

记 $\mathfrak{z}_x^{cl}=\mathfrak{z}_{\mathfrak{g}}^{cl}(\mathcal{O}_x)$ 为 1.2.2 中的分次泊松代数 $P_{(\mathfrak{g}\otimes K_x,G(\mathcal{O}_x))}=\operatorname{Sym}(\mathfrak{g}\otimes K_x/\mathcal{O}_x)^{G(\mathcal{O}_x)}$。我们将构造典范的 $\mathbb{G}_m$-等变同构 $\operatorname{Spec}\mathfrak{z}_x^{cl}\xrightarrow{\sim}\operatorname{Hitch}_x$（$\mathfrak{z}_x^{cl}$ 上的 $\mathbb{G}_m$-作用与分次所诱导的作用相反；参看 2.2.2 末尾）。

留数配对把 $(K_x/\mathcal{O}_x)^{*}$ 与 $\omega_{\mathcal{O}_x}$ 等同，故 $\operatorname{Spec}\operatorname{Sym}(\mathfrak{g}\otimes K_x/\mathcal{O}_x)=\mathfrak{g}^{*}\otimes\omega_{\mathcal{O}_x}$。投影 $\mathfrak{g}^{*}\to\mathcal{C}$ 给出仿射概形的态射 $\mathfrak{g}^{*}\otimes\omega_{\mathcal{O}_x}\to\operatorname{Hitch}_x$。它 $G(\mathcal{O}_x)$-不变，故诱导态射 $\operatorname{Spec}\mathfrak{z}_x^{cl}\to\operatorname{Hitch}_x$。为说明这是同构，须证明：$\mathfrak{g}^{*}\otimes\omega_{\mathcal{O}_x}$ 上每个 $G(\mathcal{O}_x)$-不变正则函数都来自 $\operatorname{Hitch}_x$ 上唯一的一个正则函数。显然可把 $\mathfrak{g}^{*}\otimes\omega_{\mathcal{O}_x}$ 换成 $\mathfrak{g}^{*}\otimes\mathcal{O}_x=\operatorname{Paths}(\mathfrak{g}^{*})$，把 $\operatorname{Hitch}_x$ 换成 $\operatorname{Paths}(\mathcal{C})$（对概形 $Y$，记 $\operatorname{Paths}(Y)$ 为态射 $\operatorname{Spec}\mathcal{O}_x\to Y$ 的概形）。$\mathfrak{g}^{*}$ 的正则元构成开子集 $\mathfrak{g}^{*}_{reg}$，使得 $\operatorname{codim}(\mathfrak{g}^{*}\setminus\mathfrak{g}^{*}_{reg})>1$。于是可把 $\operatorname{Paths}(\mathfrak{g}^{*})$ 换成 $\operatorname{Paths}(\mathfrak{g}^{*}_{reg})$。因态射 $\mathfrak{g}^{*}_{reg}\to\mathcal{C}$ 光滑且满，并且 $G$ 在其纤维上的作用传递，故得证。

\bdp{2.4.2} 由 1.2.2，$\mathfrak{z}_x^{cl}=P_{(\mathfrak{g}\otimes K_x,G(\mathcal{O}_x))}$ 是泊松代数。事实上，$\mathfrak{z}_x^{cl}$ 上的泊松括号为零，因为 1.2.7 中的态射 $a^{cl}:(\operatorname{Sym}(\mathfrak{g}\otimes K_x))^{\mathfrak{g}\otimes K_x}\to\mathfrak{z}_x^{cl}$ 是满射（例如可由 2.4.1 中 $\mathfrak{z}_x^{cl}$ 的描述推出），并且 $(\overline{\operatorname{Sym}}(\mathfrak{g}\otimes K_x))^{\mathfrak{g}\otimes K_x}$ 是 $\overline{\operatorname{Sym}}(\mathfrak{g}\otimes K_x)$ 的泊松中心。

\emph{注}（读者可跳过）。事实上，对任意代数群 $G$，自然态射 $\bar{a}^{cl}:(\overline{\operatorname{Sym}}(\mathfrak{g}\otimes K_x))^{G(K_x)}\to\mathfrak{z}_x^{cl}=\mathfrak{z}_{\mathfrak{g}}^{cl}(\mathcal{O}_x)$ 是满射，因此 $\mathfrak{z}_x^{cl}$ 上的泊松括号为零。下面的证明是 Feigin--Frenkel [FF92] 第 200--202 页论证的「经典极限」。把 $\mathcal{O}_x$ 与 $K_x$ 分别等同于 $\mathcal{O}:=\mathbb{C}[[t]]$ 与 $K:=\mathbb{C}((t))$。设 $f$ 是 $\mathfrak{g}^{*}\otimes\mathcal{O}$ 上的 $G(\mathcal{O})$-不变正则函数。须把它延拓为归纳概形 $\mathfrak{g}^{*}\otimes K:=\varinjlim\mathfrak{g}^{*}\otimes t^{-n}\mathcal{O}$ 上的 $G(K)$-不变正则函数 $\tilde{f}$（事实上 $\mathfrak{g}^{*}$ 可换成任意有限维 $G$-模）。至于

\SourcePage{23}

$\varphi\in\mathfrak{g}^{*}((t))$，定义 $h_{\varphi}\in\mathbb{C}((\zeta))$ 为
\[
h_{\varphi}(\zeta)=f\left(\sum_{k=0}^{N}\varphi^{(k)}(\zeta)t^{k}/k!\right),
\]
其中 $N$ 充分大（$h_{\varphi}$ 良定义，因为存在 $m$ 使得 $f$ 来自 $\mathfrak{g}^{*}\otimes(\mathcal{O}/t^{m}\mathcal{O})$ 上的函数）。把 $h_{\varphi}(\zeta)$ 写成 $\sum_n h_n(\varphi)\zeta^{n}$。函数 $h_n:\mathfrak{g}^{*}\otimes K\to\mathbb{C}$ 是 $G(K)$-不变的。置 $\tilde{f}:=h_0$。

\bdp{2.4.3} 由 2.3 与 1.2.6，有态射
\[
h_x^{cl}:\mathfrak{z}_x^{cl}\to\Gamma(\operatorname{Bun}_G,P),
\]
它类似于 1.2.3 中的态射 $h^{cl}$。为把它与 $h_X^{cl}$ 比较，考虑仿射概形的闭嵌入 $\operatorname{Hitch}(X)\hookrightarrow\operatorname{Hitch}_x$，它把 $\mathcal{C}_{\omega_X}$ 的整体截面映到其在 $x$ 的形式邻域上的限制。设 $\theta_x^{cl}:\mathfrak{z}_x^{cl}\to\mathfrak{z}^{cl}(X)$ 为相应的分次代数的满态射。容易看出
\[
h_x^{cl}=h_X^{cl}\theta_x^{cl}.
\]
因 $\mathfrak{z}_x^{cl}$ 上的泊松括号为零（见 2.4.2），且 $h_x^{cl}$ 是泊松代数态射，$\operatorname{Im}h_x^{cl}=\operatorname{Im}h_X^{cl}$ 上的泊松括号也为零。于是证明了 2.2.4(i)。

\bdsec{2.5. 量子化 I}

\bdp{2.5.1} 设 $\widetilde{\mathfrak{g}\otimes K_x}$ 为 $\mathfrak{g}\otimes K_x$ 的 Kac--Moody 中心扩张，中心为 $\mathbb{C}$，由余圈 $(u,v)\mapsto\operatorname{Res}_x c(du,v)$ 定义，$u,v\in\mathfrak{g}\otimes K_x$，其中
\begin{equation}\tag{18}
c(a,b):=-\frac{1}{2}\operatorname{Tr}(\operatorname{ad}_a\cdot\operatorname{ad}_b),\qquad a,b\in\mathfrak{g}.
\end{equation}
作为向量空间，$\widetilde{\mathfrak{g}\otimes K_x}$ 等于 $\mathfrak{g}\otimes K_x\oplus\mathbb{C}\cdot 1$。我们通过对 $g\in G(K_x)$ 指定如下自同构，来定义 $G(K_x)$ 在 $\widetilde{\mathfrak{g}\otimes K_x}$ 上的伴随作用\footnote{一旦有了李代数为 $\widetilde{\mathfrak{g}\otimes K_x}$ 的 $G(K_x)$ 的中心扩张，作用 (19) 就成为真正的伴随作用（$\widetilde{\mathfrak{g}\otimes K_x}$ 的一个在 $\mathbb{C}\cdot 1$ 与 $\mathfrak{g}\otimes K_x$ 上恒等的自同构必为恒等自同构，因为 $\operatorname{Hom}(\mathfrak{g}\otimes K_x,\mathbb{C})=0$）。}：
% [Tutor Note: 原书 p.23 末句的公式 (19) 起于 p.24，不在本片段范围内。]
